\documentclass[11pt]{article}

% Status: working draft skeleton. The canonical version of this paper is the
% project's source of truth (see Guidance_Documents/study_design.md). All
% empirical claims here should be regenerated from
% ncot_divergence_pilot.ipynb and divergence_study_outputs/.

\usepackage[margin=1in]{geometry}
\usepackage{amsmath, amssymb}
\usepackage{booktabs}
\usepackage{graphicx}
\usepackage{hyperref}
\usepackage{xcolor}
\usepackage{enumitem}

\title{Narrative Chain-of-Thought:\\
A Divergence Pilot on Ethical Reasoning Under Targeted Failure Modes}
\author{Patrick Cooper}
\date{\today}

\begin{document}
\maketitle

\begin{abstract}
We test whether prompting a large language model to narrate its reasoning
through an ethical dilemma --- as a story told from the protagonist's
perspective --- produces materially different outputs than standard
chain-of-thought (CoT). Rather than grading ethical quality, which is itself a
contested research problem, we measure two forms of \emph{divergence} between
narrative-CoT and standard-CoT outputs across five scenarios designed to
trigger known failure modes of standard CoT. Tier~1 measures structural
divergence on a six-variable rubric coded by two independent judge models from
different families, with Cohen's kappa reported per variable to bound the
rubric's noise floor. Tier~2 measures conclusion divergence as Jensen-Shannon
divergence between condition-conditional decision distributions on a fixed
per-scenario decision taxonomy, with bootstrap confidence intervals. We find
that on \texttt{gpt-5.4-nano}, narrative CoT diverges sharply from standard
CoT on the structural dimensions where standard CoT empirically fails (most
clearly on uncertainty articulation), while leaving distributions unchanged
on scenarios where standard CoT does not exhibit the targeted failure mode in
the first place. We argue this targeted-divergence pattern, not uniform
divergence, is the empirical signature that motivates further work on
narrative-structured reasoning.
\end{abstract}

\section{Introduction}

\textit{To be drafted from the author's notes. The minimum claim is: a minimal
prompt-level intervention --- asking the model to narrate rather than reason
in the abstract --- is sufficient to produce divergent structural and
conclusion outputs on scenarios that target known failure modes of standard
chain-of-thought.}

\section{Related work}

\textit{To be drafted: chain-of-thought prompting, faithfulness of CoT,
narrative reasoning in cognitive science, ethical-reasoning benchmarks for
LLMs.}

\section{Method}
\label{sec:method}

\subsection{Conditions}

Three prompting conditions, identical decoding parameters across all three:

\begin{enumerate}[leftmargin=2em]
  \item \textbf{baseline\_io}: bare input/output, ``Answer directly and concisely.''
  \item \textbf{standard\_cot}: ``Think step by step, then give your answer.''
  \item \textbf{narrative\_cot}: a five-section narrated reasoning prompt that
    asks the model to (1)~characterise the decision-maker, (2)~identify all
    stakeholders whose lives intersect with the decision, (3)~narrate the
    consequences of each available action at least two steps out, (4)~articulate
    what remains genuinely uncertain, and (5)~state the decision plainly.
    The prompt deliberately omits any formal apparatus.
\end{enumerate}

\subsection{Scenarios}

Five scenarios, surface-novel but structurally clear, each designed to
trigger one or two specific failure modes of standard CoT: scarce
experimental treatment allocation; pharmaceutical data suppression;
disagreeing siblings and dementia care; autonomous emergency vehicle policy;
research participant with changing capacity. Each scenario carries a fixed
\emph{decision taxonomy} of canonical action labels used by the Tier~2
decision extractor.

\subsection{Tier 1: structural divergence}

Per output, two independent judge models from different families
(\texttt{gpt-4o-mini} primary, \texttt{gpt-5.4-nano} secondary) code:
refusal, truncation, action commitment, stakeholder count, max causal hops
(0--5, with worked examples), uncertainty articulation (0--3), and named
ethical frameworks invoked. Reported with means by condition, Mann--Whitney
$U$ tests, Cliff's delta as the magnitude effect size, and bootstrap 95\% CIs
on Cliff's delta.

Cohen's kappa (linearly weighted for ordinal codes) is reported per
variable as the inter-judge agreement noise floor; variables with $\kappa <
0.4$ are flagged.

\subsection{Tier 2: conclusion divergence}

Each output is mapped onto its scenario's fixed decision taxonomy by an
extractor LLM. Conclusion divergence is operationalised as
\[
  \mathrm{JSD}\bigl(P_{\text{narr}} \,\big\Vert\, P_{\text{std}}\bigr)
\]
where $P_{\text{narr}}$ and $P_{\text{std}}$ are the empirical decision
distributions per scenario, JSD is in $[0, 1]$ using $\log_2$, and bootstrap
95\% CIs are computed on $N_{\text{boot}}$ resamples. We report
\[
  \text{divergence excess}
  := \mathrm{JSD}(P_{\text{narr}}, P_{\text{std}})
  - \max\bigl(\mathrm{JSD}_{\text{narr,internal}},
              \mathrm{JSD}_{\text{std,internal}}\bigr)
\]
where the within-condition JSDs are computed by random splitting; this
controls for the within-condition decision variance and reports
cross-condition divergence \emph{beyond} that noise floor.

\subsection{Targeted-failure-mode analysis}

For each (scenario, target failure mode) pair we report whether the failure
mode actually fires in standard CoT (using a per-mode operationalisation,
e.g.\ \texttt{uncertainty\_score} $\le 1$ for \texttt{uncertainty\_suppression}),
its fire rate in narrative CoT, and the corresponding Tier~2 divergence
excess. Aggregating across scenarios that target the same mode lets us ask
the email's central question: is divergence concentrated on the targeted
modes, or distributed uniformly?

\section{Results}
\label{sec:results}

\textit{To be filled in from the current notebook run. The headline numbers
in this section should be regenerated from \texttt{divergence\_study\_outputs/}
on every paper draft.}

\subsection{Tier 1: structural divergence}

\textit{Insert table \texttt{tier1\_effect\_sizes.csv} and figure
\texttt{tier1\_structural\_metrics.png} once the paper draft is final.}

\subsection{Inter-judge agreement}

\textit{Insert \texttt{interjudge\_agreement.csv}.}

\subsection{Tier 2: conclusion divergence}

\textit{Insert table \texttt{tier2\_summary.csv} and figure
\texttt{tier2\_decision\_distributions.png}.}

\subsection{Targeted failure-mode pattern}

\textit{Insert tables \texttt{failure\_mode\_profile.csv} and
\texttt{failure\_mode\_aggregate.csv}, plus figure
\texttt{divergence\_profile\_heatmap.png}.}

\section{Discussion}

\textit{To be drafted. The structure of the discussion should mirror the
interpretation guide in Section~8 of the notebook: read robustness first,
then targeted firing, then magnitude.}

\section{Limitations}

\begin{itemize}[leftmargin=2em]
  \item Two generation models, both from OpenAI. Effects observed here may
    not transfer to other vendors. The structural pattern is consistent
    across both, but the conclusion-divergence pattern differs sharply:
    \texttt{gpt-4o} is much more steerable by narrative framing than
    \texttt{gpt-5.4-nano}, suggesting the reasoning model is partly already
    doing narrative reasoning under standard CoT.
  \item Single primary judge (\texttt{claude-sonnet-4-6}) and single
    secondary judge (\texttt{gpt-4o-mini}). A third judge from a third
    vendor would harden the inter-judge agreement floor.
  \item $N=20$ samples per cell yields wide bootstrap CIs (e.g.\
    JSD$(narr, std)\in[0.07, 0.55]$ on \texttt{hospital\_allocation}). 
    Production claims need $N \ge 30$.
  \item One of the five hypothesised failure modes
    (\texttt{premature\_refusal}) does not fire on either generation model
    in this pilot, so that hypothesis is not tested.
  \item Inter-judge $\kappa$ on \texttt{max\_causal\_hops} remains low
    ($\kappa \approx 0.08$) even after rubric tightening with calibration
    anchors. Direction agreement is reported separately and survives, but
    the absolute count is unreliable.
  \item No human evaluation (Tier~3 deferred).
\end{itemize}

\section{Future work}

Listed in detail in Section~9 of the notebook. The most important next steps
are: (1)~scale to $N \ge 30$ and 20--30 scenarios; (2)~add a generation model
that exhibits the failure modes the current model does not; (3)~add a third,
cross-vendor judge; (4)~run human pairwise preference (Tier~3); (5)~ablate the
narrative-CoT prompt to identify which sub-instructions carry the
divergence.

\section*{Reproducibility}

All generation, judging, and analysis code lives in
\texttt{ncot\_divergence\_pilot.ipynb}. Per-sample artefacts (model outputs,
both judges' codings, and decision-extractor outputs) are cached in
\texttt{divergence\_study\_outputs/}. The study design is locked in
\texttt{Guidance\_Documents/study\_design.md}; this paper should be
regenerated whenever that document changes.

\end{document}
